\section{Модель и аксиоматика}

\subsection{Непротиворечивость}

Интерпретацией системы аксиом называется выбор конкретных множеств и отношений на них в качестве основных понятий. Если в интерпретации выполнены все аксиомы, она называется моделью системы~[3]. Модель Кэли~--- Клейна построена средствами евклидовой геометрии: её точки, прямые и движения~--- евклидовы точки, хорды и проективные преобразования, сохраняющие окружность. В разделах~1 и~2 проверено, что в ней выполняются аксиомы I--IV групп и аксиома Лобачевского.

Отсюда следует, что геометрия Лобачевского непротиворечива, если непротиворечива геометрия Евклида. Действительно, если бы из аксиом Лобачевского выводились утверждение и его отрицание, то, переводя вывод на язык модели, мы получили бы противоречие в евклидовой геометрии. Такая непротиворечивость называется относительной~[3,~8]: модель не доказывает, что геометрия Лобачевского свободна от противоречий сама по себе, она показывает, что эта геометрия не может оказаться противоречивой раньше, чем евклидова.

\subsection{Независимость аксиомы параллельности}

Аксиома $A$ называется независимой от системы аксиом $\Sigma$, если ни $A$, ни её отрицание не выводятся из $\Sigma$. Чтобы это доказать, достаточно предъявить две модели системы $\Sigma$: в одной $A$ выполняется, в другой нет~[3].

Для аксиомы параллельности Евклида и системы аксиом I--IV такими моделями служат евклидова плоскость и модель Кэли~--- Клейна. На евклидовой плоскости через точку вне прямой проходит ровно одна прямая, не пересекающая данную. В модели таких прямых бесконечно много (п.~2.1). Если бы аксиома параллельности выводилась из I--IV, она была бы верна в обеих моделях. Если бы выводилось её отрицание, оно тоже было бы верно в обеих. Значит, аксиома параллельности не зависит от аксиом I--IV. Именно поэтому двухтысячелетние попытки вывести V постулат из остальных аксиом не могли закончиться успехом.

\subsection{Что модель показывает и чего не доказывает}

Модель отвечает на вопрос, что означает утверждение геометрии Лобачевского, и позволяет проверить его на сколь угодно многих конкретных фигурах. Например, сумма углов треугольника, вычисленная в модели для любого положения вершин, оказывается меньше $2d$. Однако никакое число таких проверок не заменяет доказательства: модель лишь одна из интерпретаций, и рассуждение, опирающееся на свойства хорд или на вид рисунка, может использовать то, чего нет в аксиомах. Поэтому в разделах~1--3 каждое утверждение сначала доказывается средствами абсолютной геометрии или ссылкой на~[1,~2], а модель используется для интерпретации и для вычисления величин.

Есть и обратная польза. В модели легко увидеть, какое утверждение верно, прежде чем его доказывать. Если медленно передвигать конец одной из двух хорд, видно, как пересекающиеся прямые становятся параллельными, затем расходящимися, и как в этот момент появляется общий перпендикуляр, который сначала уходит к абсолюту, а затем возвращается. Такое наблюдение становится гипотезой, а доказательство затем проводится по аксиомам. В интерактивной модели рядом с каждым рисунком приведена формулировка утверждения, а вопросы к сцене предлагают сначала предсказать результат и только потом проверить его на рисунке.
